\newcommand{\bolding}{1}%

\documentclass[a4paper]{report}
\usepackage[utf8]{inputenc}
\usepackage[T1]{fontenc}
\usepackage{RJournal}
\usepackage{amsmath,amssymb,array}
\usepackage{booktabs}
\usepackage{graphicx}%
\usepackage{multirow}%
\usepackage{amsmath,amssymb,amsfonts}%
\usepackage{bm}%
\usepackage{amsthm}%
\usepackage{mathrsfs}%
\usepackage[title]{appendix}%
\usepackage{xcolor}%
\usepackage{textcomp}%
%\usepackage{manyfoot}%
\usepackage{microtype}%
\usepackage{booktabs}%
\usepackage{algorithm}%
\usepackage{algorithmicx}%
\usepackage{algpseudocode}%
\usepackage{listings}%
%\usepackage{pxfonts}%
%\usepackage[outputdir = ../]{minted}%
\usepackage{minted}%
\usepackage{xcolor}%
\usepackage{subcaption}
\usepackage{natbib}
\usepackage[font = small]{caption}
% \usepackage[colorlinks=true,linkcolor=black, citecolor=blue, urlcolor=purple]{hyperref}
\usepackage{cleveref}%

%% load any required packages FOLLOWING this line

\begin{document}

%% do not edit, for illustration only
\sectionhead{Contributed research article}
\volume{XX}
\volnumber{YY}
\year{20ZZ}
\month{AAAA}

%% replace RJtemplate with your article
\begin{article}
% \DeclareMathOperator*{\argmax}{arg\,max}
% \DeclareMathOperator*{\argmin}{arg\,min}
%\startlocaldefs
% my-commands
%\newcolumntype{C}[1]{>{\centering\arraybackslash}m{#1}}
\newcommand{\T}{\mathrm{\scriptscriptstyle T} }
\newcommand{\trans}{^\top} % transpose
%\newcommand{\tr}{\mbox{{\rm tr}}} % trace
\newcommand{\prm}{^{\prime}} % prime
\newcommand{\dprm}{^{\prime\prime}} % double prime
\newcommand{\fim}[1]{\trans{\bX}#1\bX} % Fisher's Information Matrix
\newcommand{\dd}[1]{\mfrac {\mathrm{d}}{\mathrm{d}#1}} % differentiation
\newcommand{\del}[2]{\frac{\partial#1}{\partial #2}} % partial-differentiation-1
\newcommand{\deltwo}[1]{\mfrac{\partial^2}{\partial #1^2}} % partial-differentiation-2
\newcommand\scalemath[2]{\scalebox{#1}{\mbox{\ensuremath{\displaystyle #2}}}}
\newcommand{\beginsupplement}{%
        \setcounter{equation}{0}
        \renewcommand{\theequation}{S\arabic{equation}}
        \setcounter{section}{0}
        \renewcommand{\thesection}{S\arabic{section}}
        \setcounter{table}{0}
        \renewcommand{\thetable}{S\arabic{table}}%
        \setcounter{figure}{0}
        \renewcommand{\thefigure}{S\arabic{figure}}%
     }
% stat-local defs
\newcommand{\Exp}{\mathbb{E}}
\newcommand{\Cov}{{\rm Cov}}
\newcommand{\Var}{{\rm Var}}
\newcommand{\vecc}{{\rm vec}} % vectorization operator for matrices
\newcommand{\vech}{{\rm vech}} % half-vectorization operator for symmetric matrices
\newcommand{\diag}{{\rm diag}} % diagonal matrix
\newcommand{\tildeK}{\widetilde{K}} %covariance function
\newcommand{\nablaa}{\widetilde{\nabla}} %unique derivatives -- vech
\newcommand{\Partials}{\if1\bolding{\widetilde{\bm{\partial}_\bs}}\fi\if0\bolding{\widetilde{\partial}_\bs}\fi}
\newcommand{\Sigmaa}{\widetilde{\Sigma}} %covariance after marginalizing Z
\newcommand{\given}{\, | \;} % conditional probability
\newcommand{\mstar}{m^*}
\def\A{\mathscr{A}}
\def\C{\mathscr{C}}
\def\G{\if1\bolding{\bm{\mathscr{G}}}\fi\if0\bolding{\mathscr{G}}\fi}
\def\I{\mathcal{I}}
\def\calK{\if1\bolding{\bm{\mathcal{K}}}\fi\if0\bolding{\mathcal{K}}\fi}
\def\L{\if1\bolding{\bm{\mathcal{L}}}\fi\if0\bolding{\mathcal{L}}\fi}
\def\M{\mathcal{M}}
\def\N{\mathcal{N}}
\def\P{\mathcal{P}}
\def\D{\mathcal{D}}
\def\Y{\if1\bolding{\bm{\mathcal{Y}}}\fi\if0\bolding{\mathcal{Y}}\fi}
\def\Z{\if1\bolding{\bm{\mathcal{Z}}}\fi\if0\bolding{\mathcal{Z}}\fi}
\def\R1{\mbox{$\mathfrak{R}$}}
\def\R2{\mbox{$\mathfrak{R}^2$}}
\def\R3{\mbox{$\mathfrak{R}^3$}}
\def\Rd{\mbox{$\mathfrak{R}^{d}$}}
\def\Rq{\mbox{$\mathfrak{R}^{q}$}}
% my theorems
\newtheorem{definition}{Definition}
\newtheorem{theorem}{Theorem}
\newtheorem{result}{Result}
\newtheorem{condition}{Condition}
\newtheorem{lemma}{Lemma}
\newtheorem{corollary}{Corollary}
\newtheorem{proposition}{Proposition}
\newtheorem{assumption}{Assumption}
\newtheorem{remark}{Remark}
% boldface: \usepackage{bm}
\def\bnabla{\if1\bolding{\bm{\nabla}}\fi\if0\bolding{\nabla}\fi}
\def\bpartial{\if1\bolding{\bm{\partial}}\fi\if0\bolding{\partial}\fi}
\def\0{\if1\bolding{\bm{0}}\fi\if0\bolding{0}\fi}
\def\1{\if1\bolding{\bm{1}}\fi\if0\bolding{1}\fi}
\def\ba{\if1\bolding{\bm{a}}\fi\if0\bolding{a}\fi}
\def\be{\if1\bolding{\bm{e}}\fi\if0\bolding{e}\fi}
\def\bu{\if1\bolding{\bm{u}}\fi\if0\bolding{u}\fi}
\def\bv{\if1\bolding{\bm{v}}\fi\if0\bolding{v}\fi}
\def\bw{\if1\bolding{\bm{w}}\fi\if0\bolding{w}\fi}
\def\bW{\if1\bolding{\bm{W}}\fi\if0\bolding{W}\fi}
\def\bx{\if1\bolding{\bm{x}}\fi\if0\bolding{x}\fi}
\def\bs{\if1\bolding{\bm{s}}\fi\if0\bolding{s}\fi}
\def\bDelta{\if1\bolding{\bm{\Delta}}\fi\if0\bolding{\Delta}\fi}
\def\bP{\if1\bolding{\bm{\mathrm{P}}}\fi\if0\bolding{P}\fi} % matrix
\def\bmu{\if1\bolding{\bm{\mu}}\fi\if0\bolding{\mu}\fi}
\def\btheta{\if1\bolding{\bm{\theta}}\fi\if0\bolding{\theta}\fi}
\def\bbeta{\if1\bolding{\bm{\beta}}\fi\if0\bolding{\beta}\fi}
\def\bgamma{\if1\bolding{\bm{\gamma}}\fi\if0\bolding{\gamma}\fi}
\def\bL{\if1\bolding{\bm{L}}\fi\if0\bolding{L}\fi}
\def\bV{\if1\bolding{\bm{\mathrm{V}}}\fi\if0\bolding{V}\fi} % matrix
\def\bn{\if1\bolding{\bm{n}}\fi\if0\bolding{n}\fi}
\def\m{\if1\bolding{\bm{m}}\fi\if0\bolding{m}\fi}
\def\bH{\if1\bolding{\bm{\mathrm{H}}}\fi\if0\bolding{H}\fi} % matrix
\def\bN{\if1\bolding{\bm{\mathrm{N}}}\fi\if0\bolding{N}\fi} % matrix
\def\bGamma{\if1\bolding{\bm{\Gamma}}\fi\if0\bolding{\Gamma}\fi}
\def\bF{\if1\bolding{\bm{F}}\fi\if0\bolding{F}\fi}
\def\bA{\if1\bolding{\bm{\mathrm{A}}}\fi\if0\bolding{A}\fi} % matrix
\def\bB{\if1\bolding{\bm{\mathrm{B}}}\fi\if0\bolding{B}\fi} % matrix
\def\bK{\if1\bolding{\bm{\mathrm{K}}}\fi\if0\bolding{K}\fi} % matrix
\def\bG{\if1\bolding{\bm{\mathrm{G}}}\fi\if0\bolding{G}\fi} % matrix
\def\bI{\if1\bolding{\bm{\mathrm{I}}}\fi\if0\bolding{I}\fi} % matrix
\def\bX{\if1\bolding{\bm{\mathrm{X}}}\fi\if0\bolding{X}\fi} % matrix
\def\bR{\if1\bolding{\bm{\mathrm{R}}}\fi\if0\bolding{R}\fi} % matrix
\def\bE{\if1\bolding{\bm{\mathrm{E}}}\fi\if0\bolding{E}\fi} % matrix
\def\bO{\if1\bolding{\bm{\mathrm{O}}}\fi\if0\bolding{O}\fi} % matrix
\def\bU{\if1\bolding{\bm{\mathrm{U}}}\fi\if0\bolding{U}\fi}
\def\bM{\if1\bolding{\bm{\mathrm{M}}}\fi\if0\bolding{M}\fi}
\title{Supplement to "nimblewomble: An R package for Bayesian Wombling with \texttt{nimble}"}
\author{by Aritra Halder and Sudipto Banerjee}

\maketitle

\begin{theorem}\label{th:1}
         On each rectilinear segment, $C_{t^*}=\{\bs_0+t\bu:t\in[0,t^*]\}$, where $\bu$ is a unit vector and $\bu^\perp$ is its normal, the terms in the cross-covariance matrix are
         
         \begin{equation*}
             \begin{split}
                 k_{ij}(t^*,t^*)&= (-1)^i\int_0^{t^*}\int_0^{t^*}\ba_i(t_1)^{\T}\;\bpartial^{i+j}K(\bDelta(t_1,t_2))\;\ba_j(t_2)\;||\bs'(t_1)||\;||\bs'(t_2)||\;dt_1\;dt_2,\\
                 &=(-1)^i\;t^*\int_{-t^*}^{t^*}\ba_i^{\T}\;\bpartial^{i+j}K(x\bu)\;\ba_j\;dx, \quad i,j = 1,2,
             \end{split}
         \end{equation*}
         where $\ba_1(t)=\bn(\bs(t))$ is the normal to the segment, $\ba_2(t)=\mathcal{E}_2\;\bn(\bs(t))\otimes \bn(\bs(t))$, where $\mathcal{E}_2= \left(\begin{smallmatrix}
        1 & & & \\ & 1 & 1 & \\ & & & 1
    \end{smallmatrix}\right)$ is an elimination matrix and $\bDelta(t_1,t_2)=\bs_2(t)-\bs_1(t)$.
\end{theorem}
    \begin{proof}
        Considering the parametric segment, $C_{t^*}=\{\bs_0+t\bu:t\in[0,t^*]\}$, where $\bu=(u_1,u_2)^{\T}$ is a unit vector, $||\bu||=1$, $\bu^\perp$ is its normal, $\bu^{\T}\;\bu^\perp=0$, $\ba_1(t)=\ba_1=\bu^\perp$ and $\ba_2(t)=\ba_2=\mathcal{E}_2(\bu^\perp\otimes \bu^\perp)$ are free of $t$, and $\bDelta(t_1,t_2)=(t_2-t_1)\bu$. The integrand in $k_{ij}(t^*,t^*)= (-1)^i\int_0^{t^*}\int_0^{t^*}\ba_i^{\T}\;\bpartial^{i+j}K((t_2-t_1)\bu)\;\ba_j\;dt_1\;dt_2$, depends only on $(t_2-t_1)$. Making a change of variable--define $x=t_2-t_1$, $y=t_2+t_1$, implying $\frac{x+y}{2}=t_2$ and $\frac{y-x}{2}=t_1$. The Jacobian is $\frac{1}{2}$. Hence, $0\leq y+x\leq 2t^*$ and $0\leq y-x\leq 2t^*$. This implies $0 \leq y \leq 2t^*$ and $-t^*\leq x \leq t^*$. Making the substitution above reduces, $k_{ij}(t^*,t^*)= \frac{(-1)^i}{2}\int_0^{2t^*}\int_{-t^*}^{t^*}\ba_i^{\T}\;\bpartial^{i+j}K(x\bu)\;\ba_j\;dx\;dy=(-1)^i\;t^*\int_{-t^*}^{t^*}\ba_i^{\T}\;\bpartial^{i+j}K(x\bu)\;\ba_j\;dx$. Setting $i=j=1$, produces the scenario in \cite{banerjee2006bayesian}. 
    \end{proof}

    \begin{theorem}\label{th:2}
        For the Mat\'ern kernel, with $\nu = 3/2$, the variance of the gradient based wombling measure is  $k_{11}(t^*,t^*)=2\sqrt{3}\sigma^2\phi t^*G\left(1,\sqrt{3}\phi t^*\right)$. In case $\nu=5/2$, the cross-covariance matrix for the wombling measures based on spatial gradient and curvature requires $k_{11}(t^*,t^*)=\frac{2\sqrt{5}}{3}\sigma^2\phi \;t^*\left\{G\left(1,\sqrt{5}\phi t^*\right)+G\left(2,\sqrt{5}\phi t^*\right)\right\}$, $k_{22}(t^*,t^*)=10\sqrt{5}\sigma^2\phi ^3\;t^*G(1,\sqrt{5}\phi t^*)$ and $k_{21}(t^*,t^*)=-k_{12}(t^*,t^*)=0$, where $G\left(a, \frac{x}{b}\right)=\int_0^{t^*}x^{a-1}e^{-\frac{x}{b}}dx$ is the lower incomplete Gamma function with shape parameter $a$ and scale parameter $b$. For the Gaussian kernel, 

        \begin{equation*}
            \bV_{\bGamma}(t^*)=2\sigma^2\sqrt{\pi}\phi\;t^*\left\{2\Phi\left(\sqrt{2}\phi t^*\right)-1\right\}\left(\begin{smallmatrix} 1 & 0\\0&6\phi^2 \end{smallmatrix}\right),
        \end{equation*}
        where $\Phi(\cdot)$ is the cdf for the standard Gaussian probability density.
    \end{theorem}
    \begin{proof}
        Using \Cref{th:1}, if $\nu = 3/2$, then we obtain $k_{11}(t^*,t^*)=3\sigma^2\phi ^2\;t^*\int_{-t^*}^{t^*}e^{-\sqrt{3}\phi  |x|}\;dx=2\sqrt{3}\sigma^2\phi t^*\;G(1,\sqrt{3}\phi t^*)$. If $\nu=5/2$, then $k_{11}(t^*,t^*)=\frac{5}{3}\sigma^2\phi \;t^*\int_{-t^*}^{t^*}(1+\sqrt{5}\phi |x|)\;e^{-\sqrt{5}\phi  |x|}\;dx=\frac{2\sqrt{5}}{3}\sigma^2\phi \;t^*\left\{G\left(1,\sqrt{5}\phi t^*\right)+G\left(2,\sqrt{5}\phi t^*\right)\right\}$. For terms $k_{21}$ and $k_{12}$, let us consider $a_1^{\T}\bpartial^3K(x\bu)\;a_2=\frac{25}{3}\sigma^2\phi e^{-\sqrt{5}\phi  |x|}|x|\big\{\ba_1^{\T} \bA_1 \ba_2 -\sqrt{5}\phi |x|\; \ba_1^{\T}\; \bA_2 \;\ba_2\big\}$, where 
        
        \begin{equation*}
            \bA_1 = \left(\begin{smallmatrix} 3u_1 & u_2 & u_1\\u_2 & u_1 & 3u_2 \end{smallmatrix}\right), \quad \bA_2=\left(\begin{smallmatrix} u_1^3 & u_1^2\;u_2 & u_2^2\;u_1\\ u_1^2\;u_2 & u_2^2\;u_1 & u_2^3 \end{smallmatrix}\right).
        \end{equation*} 
        The first and the second term both reduce to 0. Hence, $k_{21}(t^*,t^*)=-k_{12}(t^*,t^*)=0$. Next, $\ba_2^{\T} \bpartial^4K(x\bu)\;\ba_2 = \frac{25}{3}\sigma^2\phi ^4\;e^{-\sqrt{5}\phi  |x|}\big\{\ba_2^{\T} \bA_3\; \ba_2 - \sqrt{5}\phi \; |x|\; \ba_2^{\T} \bA_4\Big(\sqrt{5}\;\phi \;|x|+1\Big) \;\ba_2\big\}$, where 
        
        \begin{equation*}
            \bA_3=\left(\begin{smallmatrix} 3 & 0 & 1\\0 & 1 & 0\\1 & 0 & 3 \end{smallmatrix}\right), \quad \bA_4(x')=\left(\begin{smallmatrix}
        6\;u_1^2-\;x'\; u_1^4 & (3-x'\; u_1^2)\;u_1\;u_2 & 1-x'\;u_1^2\;u_2^2\\ (3-x'\; u_1^2)\;u_1\;u_2 & 1-x'\;u_1^2\;u_2^2 & (3-x'\;u_2^2)\;u_1\;u_2\\1-x'\;u_1^2\;u_2^2 & (3-x'\;u_2^2)\;u_1\;u_2 & 6\;u_2^2-\;x'\; u_2^4
    \end{smallmatrix}\right).
        \end{equation*}
        After some algebra, $\ba_2^{\T}\; \bpartial^4K(x\bu)\;\ba_2=25\sigma^2\phi ^4\;e^{-\sqrt{5}\phi  |x|}\big\{1-2\sqrt{5}\;\phi \;|x|\;\bu_1^\perp \;\bu_2^\perp\;(\bu_1\;\bu_2+\bu_1^\perp\; \bu_2^\perp)\big\}$. Observe that $\bu_1^\perp = \bu_2$ and $\bu_2^\perp=-\bu_1$, substituting we get $k_{22}(t^*,t^*)=10\sqrt{5}\sigma^2\phi ^3\;t^*G(1,\sqrt{5}\phi t^*)$. 
        
        For the Gaussian kernel,
    $k_{11}(t^*,t^*)=2\sigma^2\phi^2 \;t^*\int_{-t^*}^{t^*}e^{-\phi^2  x^2}\;dx=2\sigma^2\sqrt{\pi}\phi\;t^*\left\{2\Phi\left(\sqrt{2}\phi t^*\right)-1\right\}$. 
    Note that $\ba_1^{\T}\bpartial^3K(x\bu)\;\ba_2=4\sigma^2\phi ^4\;e^{-\phi^2  x^2}x\left\{\ba_1^{\T} \bA_1 \ba_2 -2\phi^2  x^2\; \ba_1^{\T}\; \bA_2 \;\ba_2\right\}$. Again, the first and second terms equate to 0 implying $k_{21}(t^*,t^*)=-k_{12}(t^*,t^*)=0$. Next, $\ba_2^{\T} \bpartial^4K(x\bu)\;\ba_2 = 4\sigma^2\phi^4\;e^{-\phi^2 x^2}\big\{\ba_2^{\T} \bA_3\; \ba_2 - 2\phi^2 \; x^2\; \ba_2^{\T} \bA_4(2\phi^2 x^2) \;\ba_2\big\}$. After some algebra, $\ba_2^{\T}\; \bpartial^4K(x\bu)\;\ba_2=12\sigma^2\phi^4\;e^{-\phi^2  x^2}\big\{1-4\;\phi^2 \;x^2\;\bu_1^\perp\; \bu_2^\perp\;(\bu_1\;\bu_2+\bu_1^\perp\;\bu_2^\perp)\big\}$. 
    Substituting, we get $k_{22}(t^*,t^*)=12\sigma^2\sqrt{\pi}\phi^3 t^*\left\{2\Phi\left(\sqrt{2}\phi t^*\right)-1\right\}$ resulting in the required expression for $\bV_{\bGamma}(t^*)$.
    \end{proof}
    \Cref{th:2} interestingly shows that the posited cross-covariance matrix, $\bV_{\bGamma}(t^*)$, reduces to a \emph{variance-covariance matrix}. It has a simpler form for the Gaussian case when compared to Mat\'ern with $\nu =\frac{5}{2}$. % The required expressions for $\bpartial^{i+j}K(x\bu)$, $i,j=1,2$, used for the above construction can be found in \cref{sssec:m52,sssec:gauss} (also see Section S3, Supplement of \cite{halder2024bayesian}). 
    Finally, considering the covariance between $Z(\bs_i)$, $i=1,\ldots, N$ and the wombling measures---in the Gaussian case, closed-forms are available \citep[see, e.g.,][end of Section 3]{halder_bayesian_2024}. The inferential exercise of wombling does not require quadrature when using a Gaussian kernel however, only one-dimensional quadrature is required for the same when using a Mat\'ern kernel.

    \bibliography{nimwomb}

\address{Aritra Halder\\
  Department of Biostatistics \& Epidemiology\\
  711 Nesbitt Hall, 3215 Market St., Philadelphia, PA 19104\\
  https://orcid.org/0000-0002-5139-3620\\
  \email{aritra.halder@drexel.edu}}

\address{Sudipto Banerjee\\
  Department of Biostatistics\\
  University of California, Los Angeles\\
  650 Charles E. Young Dr. S., Los Angeles, CA 90095\\
  https://orcid.org/0000-0002-2239-208X\\
  \email{sudipto@ucla.edu}}

\end{article}

\end{document}